\documentclass[a4paper,12pt]{article}

\usepackage{tgpagella}         % Palatino-like, modern, beautiful
\usepackage{microtype}         % 细节微调

\usepackage{setspace}
\setstretch{1.15}
\usepackage[margin=1in]{geometry}

\usepackage{hyperref} % for href 
\usepackage{amsmath,amsthm,amsfonts}                % AMS Math
\usepackage{thmtools}                               % Theorem Tools
\usepackage{bm}                                     % Bold Math
\usepackage{graphicx}
\usepackage{float} % can use [H] so the figure does not float
\usepackage{listings} %inline code
\usepackage{codeformat}

\usepackage[style=alphabetic,maxcitenames=1]{biblatex}
\addbibresource{reference.bib}

% Adjust the backgound of the page
\usepackage{pagecolor}
\definecolor{softivory}{RGB}{252, 249, 237}
\pagecolor{softivory}
% \nopagecolor % uncomment this when submit



%────────────────────────────────────────────────────────────────────────────────────────────────────────────────────────────────────────────────────
%────────────────────────────────────────────────────────────────────────────────────────────────────────────────────────────────────────────────────
\author{
  Author 1 \\
  Institute 1 \\
  \texttt{author1@gmail.com}
  \and
  Author 2 \\
  Institute 2 \\  
  \texttt{author2@gmail.com}
}

\date{\today}

\title{\textbf{Title of the Article}}


%────────────────────────────────────────────────────────────────────────────────────────────────────────────────────────────────────────────────────
% Statrt the content
%────────────────────────────────────────────────────────────────────────────────────────────────────────────────────────────────────────────────────
\begin{document}
\maketitle

\begin{abstract}
  This is the abstract

  \textbf{Keywords:} \textit{keyword1, keyword2}
\end{abstract}

\section{Background}
The MLIR project \cite{10.1109/CGO51591.2021.9370308} is a novel approach to building reusable 
and extensible compiler infrastructure. 
MLIR aims to address software fragmentation, 
improve compilation for heterogeneous hardware, 
significantly reduce the cost of building domain specific compilers, 
and aid in connecting existing compilers together.

\section{MLIR Code Example}
\begin{lstlisting}[style=MyMLIR]
func.func @foo(%arg0: memref<?xf32>) -> () {
  "use"(%arg0) : (memref<?xf32>) -> ()
  return
}

// Gets converted to the following
// (using type alias for brevity):
!llvm.memref_1d = !llvm.struct<(ptr, ptr, i64, array<1xi64>, array<1xi64>)>

llvm.func @foo(%arg0: !llvm.ptr,       // Allocated pointer.
               %arg1: !llvm.ptr,       // Aligned pointer.
               %arg2: i64,             // Offset.
               %arg3: i64,             // Size in dim 0.
               %arg4: i64) {           // Stride in dim 0.
  // Populate memref descriptor structure.
  %0 = llvm.mlir.undef : !llvm.memref_1d
  %1 = llvm.insertvalue %arg0, %0[0] : !llvm.memref_1d
  %2 = llvm.insertvalue %arg1, %1[1] : !llvm.memref_1d
  %3 = llvm.insertvalue %arg2, %2[2] : !llvm.memref_1d
  %4 = llvm.insertvalue %arg3, %3[3, 0] : !llvm.memref_1d
  %5 = llvm.insertvalue %arg4, %4[4, 0] : !llvm.memref_1d

  // Descriptor is now usable as a single value.
  "use"(%5) : (!llvm.memref_1d) -> ()
  llvm.return
}
\end{lstlisting}


\section{Results}

\section{Conclusion}

\section{Acknowledgment}

\pagestyle{plain}
\printbibliography

\end{document}
